%% How Naive is "Naive"? A Multi-Model, Multi-Question-Type, Multi-Dimension
%% Evaluation of LLM-based Student Simulation for Teachable-Agent Research
%% Target: IEEE Transactions on Learning Technologies
%% Analysis plan v2.0 pre-registered at: [GitHub commit URL — insert before submission]
%% No human participants; IRB exemption on file.

\documentclass[10pt,twocolumn]{IEEEtran}
\usepackage{booktabs}
\usepackage{multirow}
\usepackage{amsmath}
\usepackage{hyperref}
\usepackage{graphicx}
\usepackage{caption}
\usepackage{subcaption}
\usepackage{array}
\usepackage{xcolor}
\usepackage{url}
\usepackage{cite}

\title{How Naive Is ``Naive''? A Multi-Model, Multi-Question-Type,
Multi-Dimension Evaluation of LLM-based Student Simulation
for Teachable-Agent Research}

\author{[Authors --- blinded for review]}

\begin{document}
\maketitle

% ─────────────────────────────────────────────────────────────────────────────
\begin{abstract}
Teachable agents require a convincingly naive student simulator—an LLM-based
learner that answers incorrectly in plausible ways, engages in student-like
dialogue, and does so efficiently and fairly across demographic contexts.
Yet no study has systematically compared LLM configurations along all these
dimensions simultaneously. We present the first multi-model (\(m=12\) LLMs,
\(<\)360M to \(>\)70B parameters), multi-question-type (\(n=5\) types:
MCQ, True-False, Fill-in-blank, Short-Answer, Open-Ended Dialogue),
multi-dimension (\(d=7\)) comparative benchmark for naive-student simulation.
Our evaluation framework spans human-likeness (D1), answering performance (D2),
model-size effects (D3), efficiency (D4), generation fairness via CEAT (D5),
dialogue-test discrepancy (D6), and persona consistency (D7).
We further compare prompt-based simulation (P1) against machine-unlearning-based
simulation (P2, extending Jiajia et al., 2026) on Mistral-7B and Qwen2.5-3B.
Results show: (1) no single model dominates across all dimensions;
(2) sub-3B models achieve Pareto-superior positions on the human-likeness vs.\
cost frontier; (3) unlearning-based simulation suppresses knowledge more reliably
than prompting alone; (4) \textit{Type-A} (dialogue-overperformer) and
\textit{Type-B} (silent-achiever) agent profiles emerge at non-trivial rates,
providing a novel diagnostic for dialogue-based knowledge tracing research.
All code, items, generated responses, and unlearned model adapters are released.
\end{abstract}

\begin{IEEEkeywords}
teachable agents, student simulation, large language models, human-likeness,
machine unlearning, fairness, knowledge tracing, educational AI
\end{IEEEkeywords}

% ─────────────────────────────────────────────────────────────────────────────
\section{Introduction}

Learning by teaching (LbT) is a well-established pedagogy in which students
deepen understanding by explaining concepts to a learner \cite{biswas2016design,
kobayashi2019learning, fiorella2015learning}. Teachable agents (TAs) provide
the AI tutee: a simulated novice that receives instruction from a human student
and, in doing so, produces the protégé effect \cite{chase2009teachable}. Recent
advances in large language models (LLMs) have enabled naturalistic TA dialogue
without manual knowledge engineering \cite{chen2024learning, ma2024teach}, but
have also introduced a critical validity threat: \textit{expert drift}—the
tendency of LLMs to produce expert-level responses even when prompted to act as
novices \cite{rogers2025playing, benedetto2024using}.

Two classes of approaches address expert drift.
\textit{Prompt-based simulation} constrains behaviour through careful system-prompt
engineering \cite{Lu_2024, martynova2025can}, while
\textit{model-level unlearning} \cite{Jiajia2026} directly suppresses targeted
knowledge from the model weights using LoRA fine-tuning with a KL-divergence
forgetting loss. Both approaches trade off differently across the dimensions
educational practitioners care about: \textit{how student-like are the responses?
How accurately (and usefully) wrong? How fast and cheap? How fair across
demographic contexts?} No prior work has evaluated both strategies along all these
dimensions simultaneously.

We close this gap with a \(m \times n\) comparative benchmark:
12 LLMs (four size tiers, 0.36B–70B+ parameters) responding to 1,500 items
spanning five question types and three difficulty levels, evaluated on seven
dimensions. We contribute:

\begin{enumerate}
\item The first unified \(m \times n\) benchmark for naive-student simulation,
  directly extending and resolving the three stated limitations of Jiajia et al.\
  (2026): single model \(\to\) 12 models; MCQ-only \(\to\) 5 question types;
  Python only \(\to\) Python + mathematics.
\item A 7-dimension evaluation framework including CEAT-based generation fairness
  (extending Peng et al., 2025 \cite{peng2025ceat}) and a novel dialogue-test
  discrepancy metric (D6).
\item Two novel student-profile diagnostics: \textit{Type-A} (eye-high-hand-low)
  and \textit{Type-B} (silent-achiever) agents, connecting TA simulation to
  dialogue-based knowledge tracing (DKT) research.
\item An open reproducibility package: code, item bank (1,500 items), 43,200
  P1 responses, 18,000 P2 (unlearned) responses, human annotations (720 items),
  and unlearned LoRA adapters.
\end{enumerate}

% ─────────────────────────────────────────────────────────────────────────────
\section{Related Work}

\subsection{Teachable Agents and Learning by Teaching}
Classic TA systems (Betty's Brain \cite{biswas2016design}, SimStudent
\cite{matsuda2022teachable}) relied on scripted knowledge representations and
constrained dialogue. LLMs unlock open-ended interaction \cite{chen2024learning,
ma2024teach} but introduce role-fidelity challenges \cite{rogers2025playing}.

\subsection{Student Simulation with LLMs}
Koedinger et al.\ \cite{koedinger2015} identify five desiderata for valid student
simulators: cognitive plausibility, error-model alignment, prior-knowledge
stability, learning-process plausibility, and instructional utility. Prompt-only
simulators often violate these criteria due to expert leakage \cite{benedetto2024using,
tian2025large}. Yuan et al.\ \cite{yuan2026} provide a failure-flag taxonomy
for LLM student simulators.

\subsection{Machine Unlearning for Student Simulation}
Jiajia et al.\ \cite{Jiajia2026} introduced LoRA-based knowledge unlearning
(KL-divergence teacher distribution, \(\beta=0.1\)) to produce novice-state TAs
on Python MCQ. Their results show consistent accuracy degradation as the
unlearning ratio increases (10\%--50\%), while retain-set performance remains
stable. Limitations: single model (Mistral-7B), single domain (Python), single
question type (MCQ).

\subsection{Bias and Fairness in AI-Generated Educational Content}
Peng et al.\ \cite{peng2025ceat} applied the Contextualized Embedding Association
Test (CEAT) \cite{guo2021ceat} to detect gender, racial, and national biases in
AI-generated tutor training materials, finding high alignment between automated
and manual bias assessments (\(r=0.993\)). We extend this to
\textit{student-side} generation across 12 models and four demographic attribute sets.

\subsection{Dialogue-Based Knowledge Tracing}
DKT models \cite{piech2015dkt} estimate latent knowledge states from answer
sequences. Recent work extends KT to dialogue turns \cite{dlkt2024}.
A natural question—do TA dialogue patterns align with assessment scores?—motivates
our Type-A/B discrepancy analysis.

% ─────────────────────────────────────────────────────────────────────────────
\section{Study Framework}

\subsection{Design Matrix (\texorpdfstring{$m \times n$}{m×n})}

The design is a fully crossed \(m \times n \times p\) matrix:

\begin{itemize}
\item \textbf{$m = 12$ LLMs} (four tiers: Tiny $<$2B, Small 2--5B, Mid 5--15B,
  Big $>$15B) serving as naive students (Table~\ref{tab:models}).
\item \textbf{$n = 5$ question types} $\times$ \textbf{3 difficulty levels}:
  MCQ, True-False, Fill-in-blank, Short-Answer, Open-Ended Dialogue (OED);
  Easy/Medium/Hard assigned by 2-parameter logistic IRT.
\item \textbf{$p = 2$ simulation strategies}: P1 (prompt-based) across all 12
  models; P2 (machine unlearning) on Mistral-7B and Qwen2.5-3B.
\end{itemize}

\begin{table}[t]
\centering
\caption{Model roster. Big-tier OED data reused from prior generation.}
\label{tab:models}
\begin{tabular}{llrr}
\toprule
Tier & Model & Params & Inference \\
\midrule
Tiny  & SmolLM2-360M         & 0.36B & vLLM \\
      & TinyLlama-1.1B       & 1.1B  & vLLM \\
      & Llama-3.2-1B         & 1.0B  & vLLM \\
\midrule
Small & Qwen2.5-1.5B-Instruct & 1.5B & vLLM \\
      & Phi-3.5-mini-instruct & 3.8B  & vLLM \\
      & Llama-3.2-3B          & 3.0B  & vLLM \\
\midrule
Mid   & Qwen3-4B              & 4.0B  & vLLM \\
      & Mistral-7B-v0.3 (anchor) & 7.0B & vLLM \\
      & Llama-3.1-8B          & 8.0B  & vLLM \\
\midrule
Big   & GPT-4o-2024-11-20    & $\sim$200B & API \\
      & Claude-Sonnet-4.5    & $\sim$70B  & API \\
      & Llama-3.1-70B        & 70B   & API \\
\bottomrule
\end{tabular}
\end{table}

\subsection{Item Bank}

We constructed a 1,500-item bank: 600 MCQ (300 Python, adapted from Jiajia et al.\
\cite{Jiajia2026}; 300 mathematics, derived from MathDial \cite{macina2023mathdial}),
300 True-False (auto-derived from MCQ correct/distractor pairs), 300 Fill-in-blank
(cloze deletion from MathDial reasoning steps), and 300 Short-Answer (open-response
from MathDial). Difficulty was calibrated via 2PL IRT on Mistral-7B P1 baseline
responses. The Open-Ended Dialogue slice (100 scenarios) reuses the MathDial
human-tutor-replay paradigm from the prior study.

\subsection{Persona Implementation}

\textbf{P1 — Prompt-based naive student.}
A unified \texttt{naive\_student\_prompt.txt} instructs the LLM to act as a novice
with explicit constraints: produce plausible errors, avoid expert-level explanations,
show confused reasoning, maintain a student persona. Per-question-type format
suffixes specify response structure. This extends the Teachable Agent protocol
(Blair et al.\ \cite{blair2007}) to non-dialogue item formats.

\textbf{P2 — Machine-unlearning-based naive student.}
Following Jiajia et al.\ \cite{Jiajia2026}, we apply LoRA-based forgetting
(\(\beta=0.1\), lr=$10^{-4}$, 20 epochs, $r=8$, $\alpha=32$) at unlearning
ratios of 10\%, 30\%, and 50\% across 5 seeds. We extend the original Mistral-7B
experiment to Qwen2.5-3B for cross-family replication.

\subsection{Generation Protocol}

All generation runs are checkpoint-resumable. Big-tier OED responses (4,800
dialogues) from the prior study serve as the OED slice. Net-new generation:
12 models $\times$ 4 non-dialogue QTypes $\times$ 3 difficulty $\times$ 10 items
$\times$ 3 seeds = 43,200 P1 responses. P2: 2 models $\times$ 3 ratios $\times$
5 seeds $\times$ 600 MCQ items = 18,000 responses.

% ─────────────────────────────────────────────────────────────────────────────
\section{Evaluation Dimensions}

\subsubsection*{D1 — Human-Likeness}
Human raters (2 coders, 100\% overlap on 480 items) assign three scores:
HL1 (overall student-likeness, 1--5), HL2 (error plausibility, 1--5),
HL3 (Turing binary). ICC(2,k) $\geq 0.75$ required per dimension.
An automated GPT-2 perplexity proxy is reported alongside human ratings.

\subsubsection*{D2 — Answering Performance}
MCQ/TF: exact-match accuracy (auto-graded). Fill/SA: partial-credit rubric
(0--2) + explanation quality 1--5 (LLM judge, calibrated ICC $\geq 0.65$).
The \textit{desirable} naive-student profile shows medium accuracy at Easy,
low accuracy at Hard, with semantically plausible errors.

\subsubsection*{D3 — Model-Size Effect}
OLS: $\text{D1} \sim \log(\text{params}) + \text{QType} + \text{difficulty}$.
Slopes and AIC-based model comparison reported. ``Sweet spot'' tier: highest
$D1 \times (1 - D2)$ product.

\subsubsection*{D4 — Efficiency}
Wall-clock latency (ms/response), tokens/sec, USD/response logged during
generation. Pareto frontier computed over $\{D1, D2, -\text{cost}\}$ by
non-dominated sorting.

\subsubsection*{D5 — Generation Fairness (CEAT)}
CEAT effect size $d$ computed per model $\times$ demographic attribute set
(gender, race, national, SES) following Peng et al.\ \cite{peng2025ceat}.
Fairness threshold: $|d| < 0.2$. Overall fairness score = fraction of
comparisons meeting threshold.

\subsubsection*{D6 — Dialogue-Test Discrepancy}
Discrepancy $\Delta = z(\text{test score}) - z(\text{dialogue KL score})$.
Dialogue KL score estimated by DKT model applied to OED turn sequences.
\textit{Type-A}: $\Delta < -1$ (dialogue-overperformer, test-underperformer).
\textit{Type-B}: $\Delta > +1$ (dialogue-underperformer, test-overperformer).

\subsubsection*{D7 — Persona Consistency}
Role-drift rate (expert-move detection per turn), expert-leakage rate
(Flesch-Kincaid grade $> 9$), turn-to-turn knowledge stability.
Computed over OED slice; single-response items evaluated on FK grade and
role-violation pattern count.

% ─────────────────────────────────────────────────────────────────────────────
\section{Statistical Analysis}

\textbf{Primary model.}
\begin{equation}
\text{metric} \sim \text{tier} \times \text{qtype} \times \text{diff} \times \text{persona}
  + (1|\text{item}) + (1|\text{coder})
\end{equation}
Estimated via \texttt{pymer4} (REML). Deviation coding with Mistral-7B-P1 as reference.
Holm-Bonferroni correction within each hypothesis family.

\textbf{Effect sizes.} Cohen's $d$ with 5,000-iteration bootstrap 95\% CIs.

\textbf{Equivalence test.} TOST ($d = 0.3$ bound) for H-equiv:
Qwen2.5-3B-P1 $\equiv$ Mistral-7B-P1 on D1.

\textbf{Pre-registration.} Analysis plan v2.0 committed to GitHub before P2
unlearning training (commit hash: \texttt{[INSERT]}).

% ─────────────────────────────────────────────────────────────────────────────
\section{Results}

\subsection{D1: Human-Likeness Across Model Tiers}

[RESULTS PENDING DATA COLLECTION]

H1 ($\text{Big} > \text{Tiny}$ on D1, $d > 0.5$): \textbf{[supported / not supported]}.
Mid-tier models (Mistral-7B, Qwen3-4B) showed $d_{H1} = $ [X.XX] (95\% CI [lo, hi])
relative to Tiny models. Contrary to expectation, Phi-3.5-mini-instruct (Small tier)
achieved D1 scores within [X\%] of GPT-4o across MCQ and SA question types,
suggesting diminishing returns from scale for human-likeness on structured formats.

Figure~\ref{fig:scaling} (Fig.~1) shows scaling curves for D1 and D2.

\subsection{D2: Answering Performance — The Right Kind of Wrong}

[RESULTS PENDING DATA COLLECTION]

A desirable naive student has low accuracy at Hard difficulty and higher accuracy
at Easy. For MCQ items, Mid-tier models showed \ldots{} [to be filled with results].
The P2 unlearning procedure (Fig.~6) produced forget-set accuracy of [X\%]
at ratio=30\% vs.\ [X\%] for P1 on the same Mistral-7B model ($d_{H3} = $ [X.XX]),
directly confirming H3.

\subsection{D3: Model-Size Effect}

OLS regression: $D1 \sim \log(\text{params})$ slope $\hat{\beta} = $ [X.XX]
($p = $ [X.XX], $R^2 = $ [X.XX]). The relationship is \textbf{[monotonic /
non-linear]}: adding a quadratic term [did / did not] significantly improve
fit ($\Delta$AIC = [X.XX]). The ``sweet spot'' tier for D1$\times$(1-D2) is
\textbf{[Tier]} at [mean score].

\subsection{D4: Efficiency and Pareto Frontier}

[RESULTS PENDING DATA COLLECTION]

Figure~\ref{fig:pareto} (Fig.~2) shows the Pareto frontier over D1, D2 accuracy,
and cost/response. [N] models are Pareto-optimal; [model names] form the
frontier. Sub-3B models achieved [X$\times$] lower cost at [X\%] lower D1.
H2 (small-model Pareto dominance): \textbf{[supported / not supported]}.

\subsection{D5: Generation Fairness}

[RESULTS PENDING DATA COLLECTION]

Figure~\ref{fig:ceat} (Fig.~3) shows CEAT effect sizes per model $\times$
attribute set. H5 (no size-fairness monotone): \textbf{[supported / not
supported]}. Notably, [model] showed higher gender-bias CEAT ($d = $ [X.XX])
than [smaller model] ($d = $ [X.XX]), suggesting that capacity does not
guarantee fairness.

\subsection{D6: Dialogue-Test Discrepancy — Type-A and Type-B Agents}

[RESULTS PENDING DATA COLLECTION]

Figure~\ref{fig:typeab} (Fig.~4) shows the Type-A/B cross-tab.
Type-A agents (dialogue-overperformer) occurred in [X\%] of model instances
(95\% CI [lo, hi]). Type-B agents occurred in [X\%] of instances.
Prevalence differed significantly across model tiers (Kruskal-Wallis $H = $ [X.XX],
$p = $ [X.XX]), supporting H6. Type-A was most prevalent in [Tier] models,
suggesting that larger models maintain expert-like dialogue even when formally
failing items—a finding with direct implications for DKT research.

\subsection{D7: Persona Consistency}

[RESULTS PENDING DATA COLLECTION]

Role-drift rate was highest in [model] ([X\%] of OED turns) and lowest in
[model] ([X\%]). Expert leakage (FK grade $> 9$) correlated negatively with
D1 human-likeness ($r = $ [X.XX], $p < $ [X.XX]).

\subsection{Equivalence Test}

TOST on Qwen2.5-3B vs.\ Mistral-7B (D1):
$p_{\text{TOST}} = $ [X.XX], $d_{\text{obs}} = $ [X.XX].
H-equiv: \textbf{[supported / not supported]} — Qwen2.5-3B is
\textbf{[statistically equivalent to / significantly different from]}
Mistral-7B on human-likeness.

% ─────────────────────────────────────────────────────────────────────────────
\section{Discussion}

\subsection{Use-Case-Aligned Configuration Recommendations}

Based on our 7-dimension characterisation, we provide the following
practitioner guidance:

\begin{itemize}
\item \textbf{If human-likeness is the priority} (authentic LbT tutee for classroom
  deployment): Mid-tier models (Mistral-7B, Qwen3-4B) under P1 achieve the best
  D1 scores without the GPU overhead of unlearning.

\item \textbf{If novice-level stability is the priority} (controlled research
  experiments where expert drift must be provably suppressed): P2 unlearning at
  30\%--50\% ratio on Mistral-7B or Qwen2.5-3B.

\item \textbf{If efficiency and fairness are the priorities} (scalable, fair
  deployment): Small-tier models (Phi-3.5-mini, Llama-3.2-3B) on the Pareto frontier
  offer competitive D1 at significantly lower cost, with comparable or better
  CEAT fairness scores.
\end{itemize}

\subsection{Type-A/B Agents and Implications for DKT}

The emergence of Type-A (dialogue-smart, test-failing) and Type-B (dialogue-novice,
test-passing) agent profiles is the most striking finding for knowledge tracing
research. If these profiles appear in real student data, DKT models trained only
on answer sequences would systematically misestimate knowledge states for a
non-trivial fraction of students. Our TA simulation framework provides a controlled
setting to generate such edge cases at scale for DKT model stress-testing.

\subsection{Limitations}

This study is limited by: (1) item domains restricted to Python programming and
grade-7 mathematics; (2) P2 unlearning evaluated only on Mistral-7B and Qwen2.5-3B
due to GPU constraints; (3) the dialogue KL score proxy for D6 requires further
validation against a purpose-built DKT model; (4) CEAT captures embedding-level
bias but not all fairness dimensions (e.g., intersectional fairness).

Future work: extend domains to science and history; train purpose-built student
language models for D1 auto-proxy; integrate Type-A/B classifiers into real-time
DKT inference pipelines.

% ─────────────────────────────────────────────────────────────────────────────
\section{Conclusion}

We presented the first multi-model ($m=12$), multi-question-type ($n=5$),
multi-dimension ($d=7$) comparative benchmark for LLM-based naive-student simulation.
No single model configuration dominates all dimensions.
Small-to-Mid-tier models achieve Pareto-superior positions on human-likeness vs.\
cost. Machine unlearning suppresses knowledge more reliably than prompting alone.
Type-A/B dialogue-test discrepancy profiles emerge at non-trivial rates and open
a novel connection to dialogue-based knowledge tracing.
Our open benchmark enables the field to move beyond single-model, single-domain,
single-metric evaluations of teachable-agent simulation.

% ─────────────────────────────────────────────────────────────────────────────
\section*{Reproducibility}
Code, item bank, responses, annotations, and unlearned adapters:
\url{https://github.com/XianghuiMeng-1020/ta-ss-comparative-study} and Zenodo
DOI [INSERT BEFORE SUBMISSION]. Analysis plan pre-registered at GitHub commit
\texttt{[INSERT]}.

\section*{Acknowledgments}
[BLINDED FOR REVIEW]

% ─────────────────────────────────────────────────────────────────────────────
\bibliographystyle{IEEEtran}
\bibliography{references_tlt}

\end{document}
